%%%%%%%%%%%%%%%%
%%% deut 2 %%%
%%%%%%%%%%%%%%%%

\Chapter{2}

\Verse{1}
{
	\hDtr1{וַנֵּפֶן וַנִּסַּע הַמִּדְבָּרָה דֶּרֶךְ יַם סוּף כַּאֲשֶׁר דִּבֶּר יְהֹוָה אֵלָי וַנָּסב אֶת הַר שֵׂעִיר יָמִים רַבִּים׃}
}
{
	\eDtr1{
		2 and we turned and traveled to the wilderness by way of
		the Red Sea as YHWH spoke to me, and we stayed around
		Mount Seir many days.
	}
}
{
%	\footnote{
%		by way of the Red Sea. The Israelites never go back near any place that
%		could be the site of the splitting of the sea. This reference therefore
%		must be to the eastern arm of the Red Sea. It is further confirmation
%		that the name of that entire body of water, Hebrew yam sûp, in the Torah
%		refers to what is known today as the Red Sea, and not to any imagined
%		“Reed Sea.” See the comments on Exod 13:8 and Num 21:4.
%	}
}

\Verse{2}
{
	\hDtr1{וַיֹּאמֶר יְהֹוָה אֵלַי לֵאמֹר׃}
}
{
	\eDtr1{“And YHWH said to me, saying,}
}
{
}

\Verse{3}
{
	\hDtr1{רַב לָכֶם סֹב אֶת הָהָר הַזֶּה פְּנוּ לָכֶם צָפֹנָה׃}
}
{
	\eDtr1{
		‘You’ve had enough of staying around this mountain. Turn
		north.
	}
}
{
}

\Verse{4}
{
	\hDtr1{וְאֶת הָעָם צַו לֵאמֹר אַתֶּם עֹבְרִים בִּגְבוּל אֲחֵיכֶם בְּנֵי עֵשָׂו הַיֹּשְׁבִים בְּשֵׂעִיר וְיִירְאוּ מִכֶּם וְנִשְׁמַרְתֶּם מְאֹד׃}
}
{
	\eDtr1{
		And command the people, saying: You’re crossing the border
		of your brothers, the children of Esau, who live in Seir.
		And they’ll be afraid of you, so be very watchful.
	}
}
{
%	\footnote{
%		they’ll be afraid of you. The shoe is on the other foot. When Jacob
%		approached his brother Esau at Seir, “Jacob was very afraid” (Gen
%		32:8,12). Now Jacob’s children are close to Seir, and Esau’s children
%		will be afraid. But YHWH tells Israel to leave them alone. The stated
%		reason is that God gave this land to Esau (Num 2:5). Another reason may
%		be that Esau did not harm Jacob and his children, and so, on the merit
%		of Esau’s act of forgiveness, Jacob’s children should now not harm
%		Esau’s children. In any case, we should take note of the benign attitude
%		that the Torah takes toward Edom, who are identified as the descendants
%		of Esau. This changes radically later on. First Edom is conquered by
%		David and then held in subjugation by Judah for many years during the
%		monarchy. The bad relations with Edom culminate following the
%		destruction of Jerusalem, with bitter denunciations of Edom (in Psalm
%		137, the book of Lamentations, and the book of Obadiah).
%	}
}

\Verse{5}
{
	\hDtr1{אַל תִּתְגָּרוּ בָם כִּי לֹא אֶתֵּן לָכֶם מֵאַרְצָם עַד מִדְרַךְ כַּף רָגֶל כִּי יְרֻשָּׁה לְעֵשָׂו נָתַתִּי אֶת הַר שֵׂעִיר׃}
}
{
	\eDtr1{
		Don’t get agitated at them, because I won’t give you any
		of their land, even as much as a footstep, because I’ve
		given Mount Seir to Esau as a possession.
	}
}
{
}

\Verse{6}
{
	\hDtr1{אֹכֶל תִּשְׁבְּרוּ מֵאִתָּם בַּכֶּסֶף וַאֲכַלְתֶּם וְגַם מַיִם תִּכְרוּ מֵאִתָּם בַּכֶּסֶף וּשְׁתִיתֶם׃}
}
{
	\eDtr1{
		Food: you shall buy from them with money, and then you’ll
		eat. And water, too: you shall purchase from them with
		money, and then you’ll drink.
	}
}
{
}

\Verse{7}
{
	\hDtr1{כִּי יְהֹוָה אֱלֹהֶיךָ בֵּרַכְךָ בְּכֹל מַעֲשֵׂה יָדֶךָ יָדַע לֶכְתְּךָ אֶת הַמִּדְבָּר הַגָּדֹל הַזֶּה זֶה אַרְבָּעִים שָׁנָה יְהֹוָה אֱלֹהֶיךָ עִמָּךְ לֹא חָסַרְתָּ דָּבָר׃}
}
{
	\eDtr1{
		Because YHWH, your God, has blessed you in all your hand’s
		work. He has known your walking in this big wilderness.
		These forty years, YHWH, your God, has been with you. You
		haven’t lacked a thing.’
	}
}
{
%	\footnote{
%		These forty years. He has known his people for forty years, and the text
%		conveys a depth of acquaintance with them that is not apparent in the
%		preceding books. Never hesitant to criticize the people of Israel, he
%		nonetheless speaks to them now with affection. He begins and ends his
%		speech with blessings on the people. In his opening he says to them
%		(1:11): May YHWH, your fathers’ God, add on to you a thousand times more
%		like you! And may He bless you as He spoke to you. And in his very last
%		words, he sings to them (33:29): Happy are you, Israel! Who is like you,
%		a people saved by YHWH. He has also known his God for forty years, but
%		their singular acquaintance—unique among all divine-human encounters in
%		the Hebrew Bible—is now among the most thunderous of the Bible’s loud
%		silences. God speaks very little to Moses in Deuteronomy, and Moses
%		never says a word to God, even in the chapters following his speech. (He
%		only quotes previous conversations he has had with God.)
%	}
}

\Verse{8}
{
	\hDtr1{וַנַּעֲבֹר מֵאֵת אַחֵינוּ בְנֵי עֵשָׂו הַיֹּשְׁבִים בְּשֵׂעִיר מִדֶּרֶךְ הָעֲרָבָה מֵאֵילַת וּמֵעֶצְיֹן גָּבֶר {ס} וַנֵּפֶן וַנַּעֲבֹר דֶּרֶךְ מִדְבַּר מוֹאָב׃}
}
{
	\eDtr1{
		So we passed by our brothers, the children of Esau, who
		live in Seir, from the way of the plain, from Eilat, and
		from Ezion-geber, and we turned and passed by way of the
		wilderness of Moab.
	}
}
{
}

\Verse{9}
{
	\hDtr1{וַיֹּאמֶר יְהֹוָה אֵלַי אַל תָּצַר אֶת מוֹאָב וְאַל תִּתְגָּר בָּם מִלְחָמָה כִּי לֹא אֶתֵּן לְךָ מֵאַרְצוֹ יְרֻשָּׁה כִּי לִבְנֵי לוֹט נָתַתִּי אֶת עָר יְרֻשָּׁה׃}
}
{
	\eDtr1{
		“And YHWH said to me, ‘Don’t oppose Moab and don’t agitate
		them with war, because I won’t give you a possession from
		its land, because I’ve given Ar as a possession to the
		children of Lot.
	}
}
{
%	\footnote{
%		Ar. A Moabite city mentioned in Num 21:15.
%	}
}

\Verse{10}
{
	\hDtr1{הָאֵמִים לְפָנִים יָשְׁבוּ בָהּ עַם גָּדוֹל וְרַב וָרָם כָּעֲנָקִים׃}
}
{
	\eDtr1{
		The Emim had lived there before, a people that was big and
		numerous and tall as giants.
	}
}
{
}

\Verse{11}
{
	\hDtr1{רְפָאִים יֵחָשְׁבוּ אַף הֵם כָּעֲנָקִים וְהַמֹּאָבִים יִקְרְאוּ לָהֶם אֵמִים׃}
}
{
	\eDtr1{
		They were also thought to be Rephaim, like the giants, but
		the Moabites call them Emim.
	}
}
{
%	\footnote{
%		Rephaim. In some places we have references to the Rephaim (and the Emim,
%		who are connected with the Rephaim here) among the indigenous peoples of
%		Canaan (Gen 14:5; Deut 2:20), but in other places the Rephaim are the
%		dead, continuing some sort of existence in an underworld (Isa 14:9;
%		26:14; Ps 88:11; see also Prov 2:18; 9:18; 21:16; Job 26:5).
%		Fifth-century Phoenician inscriptions identify the Rephaim as those whom
%		the living join in dying (KAI 13:7–8, 14:8). The term is also found in
%		Ugaritic (rapi!uma, KTU 1.61), referring to a line of dead kings and
%		heroes (cf. Isa 14:9). The deceased are invoked to assist in bestowing
%		blessings upon the reigning king. Alan Cooper has traced the etymology
%		of Rephaim to Ugaritic Rp!u, a chthonic deity and patron god of the king
%		of Ugarit, associated with healing, in the sense of granting health,
%		strength, fertility, and fecundity; hence the Hebrew r[image]p[image]’,
%		to heal. These associations are important in discussing ancestor
%		veneration in the ancient Near East, since the purposes for revering
%		one’s dead ancestors were often requests for health, strength, and
%		progeny.
%	}
}

\Verse{12}
{
	\hDtr1{וּבְשֵׂעִיר יָשְׁבוּ הַחֹרִים לְפָנִים וּבְנֵי עֵשָׂו יִירָשׁוּם וַיַּשְׁמִידוּם מִפְּנֵיהֶם וַיֵּשְׁבוּ תַּחְתָּם כַּאֲשֶׁר עָשָׂה יִשְׂרָאֵל לְאֶרֶץ יְרֻשָּׁתוֹ אֲשֶׁר נָתַן יְהֹוָה לָהֶם׃}
}
{
	\eDtr1{
		And the Horites lived in Seir before, but the children of
		Esau dispossessed them and destroyed them in front of them
		and lived in their place, as Israel did to the land of its
		possession that YHWH gave them.
	}
}
{
}

\Verse{13}
{
	\hDtr1{עַתָּה קֻמוּ וְעִבְרוּ לָכֶם אֶת נַחַל זָרֶד וַנַּעֲבֹר אֶת נַחַל זָרֶד׃}
}
{
	\eDtr1{
		Now get up and cross the Wadi Zered.’ And we crossed the
		Wadi Zered.
	}
}
{
}

\Verse{14}
{
	\hDtr1{וְהַיָּמִים אֲשֶׁר הָלַכְנוּ מִקָּדֵשׁ בַּרְנֵעַ עַד אֲשֶׁר עָבַרְנוּ אֶת נַחַל זֶרֶד שְׁלֹשִׁים וּשְׁמֹנֶה שָׁנָה עַד תֹּם כּל הַדּוֹר אַנְשֵׁי הַמִּלְחָמָה מִקֶּרֶב הַמַּחֲנֶה כַּאֲשֶׁר נִשְׁבַּע יְהֹוָה לָהֶם׃}
}
{
	\eDtr1{
		And the days that we went from Kadesh-barnea until we
		crossed the Wadi Zered were thirty-eight years, until the
		end of all the generation, the men of war, from among the
		camp, as YHWH swore to them.
	}
}
{
}

\Verse{15}
{
	\hDtr1{וְגַם יַד יְהֹוָה הָיְתָה בָּם לְהֻמָּם מִקֶּרֶב הַמַּחֲנֶה עַד תֻּמָּם׃}
}
{
	\eDtr1{
		And YHWH’s hand was against them, to eliminate them from
		among the camp until their end.
	}
}
{
}

\Verse{16}
{
	\hDtr1{וַיְהִי כַאֲשֶׁר תַּמּוּ כּל אַנְשֵׁי הַמִּלְחָמָה לָמוּת מִקֶּרֶב הָעָם׃}
}
{
	\eDtr1{
		“And it was when all the men of war had ended, dying from
		among the people,
	}
}
{
}

\Verse{17}
{
	\hDtr1{וַיְדַבֵּר יְהֹוָה אֵלַי לֵאמֹר׃}
}
{
	\eDtr1{and YHWH spoke to me, saying,}
}
{
}

\Verse{18}
{
	\hDtr1{אַתָּה עֹבֵר הַיּוֹם אֶת גְּבוּל מוֹאָב אֶת עָר׃}
}
{
	\eDtr1{‘You’re crossing the border of Moab, Ar, today.}
}
{
}

\Verse{19}
{
	\hDtr1{וְקָרַבְתָּ מוּל בְּנֵי עַמּוֹן אַל תְּצֻרֵם וְאַל תִּתְגָּר בָּם כִּי לֹא אֶתֵּן מֵאֶרֶץ בְּנֵי עַמּוֹן לְךָ יְרֻשָּׁה כִּי לִבְנֵי לוֹט נְתַתִּיהָ יְרֻשָּׁה׃}
}
{
	\eDtr1{
		And when you come close opposite the children of Ammon,
		don’t oppose them and don’t get agitated at them, because
		I won’t give you a possession from the land of the
		children of Ammon, because I’ve given it as a possession
		to the children of Lot.’
	}
}
{
%	\footnote{
%		the children of Lot. The stories in Genesis now return to mind with yet
%		another layer of significance. Israel is not to be hostile to Edom
%		(because they are Esau’s children) or Moab and Ammon (because they are
%		Lot’s children). So the stories of individuals and families in Genesis
%		now become the stories of nations in Deuteronomy. Ancient Israel
%		understood its neighboring peoples to be relatives. Israel was expected
%		to act out of kinship to them. And hostility from any of them was
%		regarded as betrayal by a family member. This attitude continues past
%		the Torah into the narrative of Israel’s history in its land.
%	}
}

\Verse{20}
{
	\hDtr1{אֶרֶץ רְפָאִים תֵּחָשֵׁב אַף הִוא רְפָאִים יָשְׁבוּ בָהּ לְפָנִים וְהָעַמֹּנִים יִקְרְאוּ לָהֶם זַמְזֻמִּים׃}
}
{
	\eDtr1{
		(It, too, was thought to be a land of Rephaim: Rephaim had
		lived there before, and the Ammonites call them Zamzummim,
	}
}
{
}

\Verse{21}
{
	\hDtr1{עַם גָּדוֹל וְרַב וָרָם כָּעֲנָקִים וַיַּשְׁמִידֵם יְהֹוָה מִפְּנֵיהֶם וַיִּירָשֻׁם וַיֵּשְׁבוּ תַחְתָּם׃}
}
{
	\eDtr1{
		a people that was big and numerous and tall as giants, but
		YHWH destroyed them in front of them, and they
		dispossessed them and lived in their place,
	}
}
{
}

\Verse{22}
{
	\hDtr1{כַּאֲשֶׁר עָשָׂה לִבְנֵי עֵשָׂו הַיֹּשְׁבִים בְּשֵׂעִיר אֲשֶׁר הִשְׁמִיד אֶת הַחֹרִי מִפְּנֵיהֶם וַיִּירָשֻׁם וַיֵּשְׁבוּ תַחְתָּם עַד הַיּוֹם הַזֶּה׃}
}
{
	\eDtr1{
		as He did for the children of Esau, who live in Seir, in
		that He destroyed the Horite in front of them, and they
		dispossessed them and have lived in their place to this
		day.
	}
}
{
}

\Verse{23}
{
	\hDtr1{וְהָעַוִּים הַיֹּשְׁבִים בַּחֲצֵרִים עַד עַזָּה כַּפְתֹּרִים הַיֹּצְאִים מִכַּפְתֹּר הִשְׁמִידֻם וַיֵּשְׁבוּ תַחְתָּם׃}
}
{
	\eDtr1{
		So the Avvim, who lived in villages as far as Gaza: the
		Caphtorites who left Caphtor destroyed them, and they
		dwelled in their place.)
	}
}
{
%	\footnote{
%		Caphtor. Crete or another of the Greek islands. It is associated with
%		the Philistines (Gen 10:14; Amos 9:7; Jer 47:4; 1 Chr 1:12).
%	}
}

\Verse{24}
{
	\hDtr1{קוּמוּ סְּעוּ וְעִבְרוּ אֶת נַחַל אַרְנֹן רְאֵה נָתַתִּי בְיָדְךָ אֶת סִיחֹן מֶלֶךְ חֶשְׁבּוֹן הָאֱמֹרִי וְאֶת אַרְצוֹ הָחֵל רָשׁ וְהִתְגָּר בּוֹ מִלְחָמָה׃}
}
{
	\eDtr1{
		‘Get up, travel, and cross the Wadi Arnon. See: I’ve put
		Sihon, king of Heshbon, the Amorite, and his land in your
		hand. Begin. Dispossess! And you shall agitate him with
		war.
	}
}
{
}

\Verse{25}
{
	\hDtr1{הַיּוֹם הַזֶּה אָחֵל תֵּת פַּחְדְּךָ וְיִרְאָתְךָ עַל פְּנֵי הָעַמִּים תַּחַת כּל הַשָּׁמָיִם אֲשֶׁר יִשְׁמְעוּן שִׁמְעֲךָ וְרָגְזוּ וְחָלוּ מִפָּנֶיךָ׃}
}
{
	\eDtr1{
		This day I’ll begin to put an awe of you and a fear of you
		on the faces of the peoples under all the skies so that
		when they’ll hear the news of you they’ll tremble and
		writhe in front of you.’
	}
}
{
}

\Verse{26}
{
	\hDtr1{וָאֶשְׁלַח מַלְאָכִים מִמִּדְבַּר קְדֵמוֹת אֶל סִיחוֹן מֶלֶךְ חֶשְׁבּוֹן דִּבְרֵי שָׁלוֹם לֵאמֹר׃}
}
{
	\eDtr1{
		“And I sent messengers from the wilderness of Kedemoth to
		Sihon, king of Heshbon, words of peace, saying,
	}
}
{
}

\Verse{27}
{
	\hDtr1{אֶעְבְּרָה בְאַרְצֶךָ בַּדֶּרֶךְ בַּדֶּרֶךְ אֵלֵךְ לֹא אָסוּר יָמִין וּשְׂמֹאול׃}
}
{
	\eDtr1{
		‘Let me pass through your land. I’ll go on the road, on
		the road! I won’t turn right or left.
	}
}
{
}

\Verse{28}
{
	\hDtr1{אֹכֶל בַּכֶּסֶף תַּשְׁבִּרֵנִי וְאָכַלְתִּי וּמַיִם בַּכֶּסֶף תִּתֶּן לִי וְשָׁתִיתִי רַק אֶעְבְּרָה בְרַגְלָי׃}
}
{
	\eDtr1{
		Food: you’ll sell me for money, and I’ll eat. And water:
		you’ll give me for money, and I’ll drink. Only let me pass
		on my feet
	}
}
{
}

\Verse{29}
{
	\hDtr1{כַּאֲשֶׁר עָשׂוּ לִי בְּנֵי עֵשָׂו הַיֹּשְׁבִים בְּשֵׂעִיר וְהַמּוֹאָבִים הַיֹּשְׁבִים בְּעָר עַד אֲשֶׁר אֶעֱבֹר אֶת הַיַּרְדֵּן אֶל הָאָרֶץ אֲשֶׁר יְהֹוָה אֱלֹהֵינוּ נֹתֵן לָנוּ׃}
}
{
	\eDtr1{
		as the children of Esau, who live in Seir, and the
		Moabites who live in Ar did for me, until I’ll cross the
		Jordan to the land that YHWH, our God, is giving us.’
	}
}
{
}

\Verse{30}
{
	\hDtr1{וְלֹא אָבָה סִיחֹן מֶלֶךְ חֶשְׁבּוֹן הַעֲבִרֵנוּ בּוֹ כִּי הִקְשָׁה יְהֹוָה אֱלֹהֶיךָ אֶת רוּחוֹ וְאִמֵּץ אֶת לְבָבוֹ לְמַעַן תִּתּוֹ בְיָדְךָ כַּיּוֹם הַזֶּה׃}
}
{
	\eDtr1{
		But Sihon, king of Heshbon, was not willing to let us pass
		through it, because YHWH, your God, had hardened his
		spirit and made his heart bold in order to put him in your
		hand—as it is this day.
	}
}
{
}

\Verse{31}
{
	\hDtr1{וַיֹּאמֶר יְהֹוָה אֵלַי רְאֵה הַחִלֹּתִי תֵּת לְפָנֶיךָ אֶת סִיחֹן וְאֶת אַרְצוֹ הָחֵל רָשׁ לָרֶשֶׁת אֶת אַרְצוֹ׃}
}
{
	\eDtr1{
		“And YHWH said to me, ‘See: I’ve begun to put Sihon and
		his land before you. Begin. Dispossess, so as to possess
		his land.’
	}
}
{
}

\Verse{32}
{
	\hDtr1{וַיֵּצֵא סִיחֹן לִקְרָאתֵנוּ הוּא וְכל עַמּוֹ לַמִּלְחָמָה יָהְצָה׃}
}
{
	\eDtr1{
		“And Sihon came out at us, he and all his people, for war
		at Jahaz.
	}
}
{
}

\Verse{33}
{
	\hDtr1{וַיִּתְּנֵהוּ יְהֹוָה אֱלֹהֵינוּ לְפָנֵינוּ וַנַּךְ אֹתוֹ וְאֶת בָּנָו וְאֶת כּל עַמּוֹ׃}
}
{
	\eDtr1{
		And YHWH, our God, put him before us, and we struck him
		and his sons and all his people.
	}
}
{
}

\Verse{34}
{
	\hDtr1{וַנִּלְכֹּד אֶת כּל עָרָיו בָּעֵת הַהִוא וַנַּחֲרֵם אֶת כּל עִיר מְתִם וְהַנָּשִׁים וְהַטָּף לֹא הִשְׁאַרְנוּ שָׂרִיד׃}
}
{
	\eDtr1{
		And we captured all his cities at that time, and we put
		every city to complete destruction: men and women and
		infants. We didn’t leave a remnant.
	}
}
{
%	\footnote{
%		complete destruction. In contexts that do not have to do with war, the
%		Hebrew word [image]erem refers to something that is devoted to God (Lev
%		27:21,28–29; Num 18:14). In contexts of war, as in this verse, øerem
%		refers to the rule, in divinely commanded wars only, against taking
%		spoils or slaves, but rather destroying all of these and thus dedicating
%		them to the deity. The point: the war is not for profit.
%	}
}

\Verse{35}
{
	\hDtr1{רַק הַבְּהֵמָה בָּזַזְנוּ לָנוּ וּשְׁלַל הֶעָרִים אֲשֶׁר לָכָדְנוּ׃}
}
{
	\eDtr1{
		Only the animals we despoiled for ourselves, and the spoil
		of the cities that we captured.
	}
}
{
}

\Verse{36}
{
	\hDtr1{מֵעֲרֹעֵר אֲשֶׁר עַל שְׂפַת נַחַל אַרְנֹן וְהָעִיר אֲשֶׁר בַּנַּחַל וְעַד הַגִּלְעָד לֹא הָיְתָה קִרְיָה אֲשֶׁר שָׂגְבָה מִמֶּנּוּ אֶת הַכֹּל נָתַן יְהֹוָה אֱלֹהֵינוּ לְפָנֵינוּ׃}
}
{
	\eDtr1{
		From Aroer, which is on the edge of the Wadi Arnon and the
		city that is in the wadi, to Gilead, there wasn’t a town
		that was too high for us. YHWH, our God, put it all before
		us.
	}
}
{
}

\Verse{37}
{
	\hDtr1{רַק אֶל אֶרֶץ בְּנֵי עַמּוֹן לֹא קָרָבְתָּ כּל יַד נַחַל יַבֹּק וְעָרֵי הָהָר וְכֹל אֲשֶׁר צִוָּה יְהֹוָה אֱלֹהֵינוּ׃}
}
{
	\eDtr1{
		Only to the land of the children of Ammon did you not come
		close, all along the Wadi Jabbok, and the cities of the
		hill country, and whatever YHWH, our God, commanded.
	}
}
{
}
